
        \documentclass[fleqn]{article}      
        \usepackage{amsmath}
        \usepackage{fontspec}
        \usepackage{kotex} % 한국어 지원  

        \begin{document}      
        쉬움: \\\\  
1. 정적분의 기호는 무엇인가요?  
   a) ∫  
   b) ∑  
   c) ∏  
   d) √ \\\\  
2. 다음 중 정적분의 정의에 해당하는 것은?  
   a) 미분  
   b) 넓이  
   c) 기울기  
   d) 평균 \\\\  
3. 다음 중 정적분의 활용 예시로 올바른 것은?  
   a) 속도 계산  
   b) 면적 계산  
   c) 기울기 계산  
   d) 최대값 찾기 \\\\  
4. 함수 f(x) = x^2의 x = 0에서 x = 2까지의 정적분은 무엇인가요?  
   a) 1  
   b) 2  
   c) 4  
   d) 8 \\\\  
5. ∫(2x + 3)dx의 결과는?  
   a) x^2 + 3x + C  
   b) x^2 + 3x  
   c) 2x^2 + 3x + C  
   d) 2x^2 + 3 + C \\\\

중간: \\\\  
1. 정적분의 기본 정리는 무엇인가요?  
   a) 미분과 적분은 서로 반대이다.  
   b) 정적분의 결과는 항상 양수이다.  
   c) 정적분은 미분의 역이다.  
   d) 정적분은 항상 상수이다. \\\\  
2. 함수 f(x) = 3x^2의 x = 1에서 x = 3까지의 정적분을 구하세요. \\\\  
3. ∫(sin x)dx의 결과는 무엇인가요?  
   a) -cos x + C  
   b) cos x + C  
   c) sin x + C  
   d) -sin x + C \\\\  
4. 정적분을 통해 구한 면적이 음수일 수 있는 경우는 무엇인가요? \\\\  
5. f(x) = x의 0에서 1까지의 정적분을 구하세요. \\\\  

어려움: \\\\  
1. 함수 f(x) = x^3의 1에서 2까지의 정적분을 구하세요. \\\\  
2. ∫(e^x)dx의 결과는 무엇인가요?  
   a) e^x + C  
   b) e^x  
   c) xe^x + C  
   d) xe + C \\\\  
3. 정적분의 기하학적 의미는 무엇인가요? \\\\  
4. 주어진 함수 f(x) = x^2 - 4x + 4의 정적분을 0에서 4까지 구하세요. \\\\  
5. ∫(1/x)dx의 결과는 무엇인가요?  
   a) ln|x| + C  
   b) 1/x + C  
   c) x + C  
   d) x^2 + C \\\\  

      
        \end{document} 
    